\section{Complexity Analysis of Deterministic DCA}
\label{sec:complexity-analysis-DCA}

This section establishes the iteration complexity of DCA-type methods for
reaching a $(\delta_\phi,\delta_g,\varepsilon)$-critical point of
$F=\phi-g$, without imposing smoothness on $g$. We analyze three
algorithmic templates against this single certificate: (1) exact and inexact
DCA, which solve the strongly convex subproblem to full or controlled accuracy;
(2) proximal DCA, which replaces that solve by a single proximal-subgradient
step; and (3) concave-side smoothing.
% , which, rather than a new method, is
shown to reproduce the same certificate the other two attain. 
% Two devices
% recur throughout, a strong-convexity reformulation that is without loss
% of generality (\Cref{ass:strong-convexity-dc}) and a Fenchel-gap toolkit
% that turns each subproblem's optimality into a transported tolerance. 

\begin{algorithm}[tbp]
    \caption{\textbf{(In)exact DCA}}
    \label{alg: DCA}
    \begin{algorithmic}[1]
        \STATE {\bfseries Input:} Initial point $\bx^0\in\dom \phi $,
        inner solver tolerance $\{\epsilon_k\}_{k=0}^{K-1}$,
        stepsize $\eta$.
        \FOR{$k=0,1,\ldots,K-1$}
            \STATE Compute $\bv^k \in \partial g(\bx^k)$.
            \STATE 
            Solve inner strongly convex subproblem:

            \STATE $\triangleright$ \textbf{Option I:} exact subproblem solve
            \begin{align}
            \label{eq:dca-exact-step}
            \bx^{k+1} = \argmin\limits_{\bx \in \RR^n } \{ \phi (\bx)- \langle \bv^k, \bx \rangle \} 
            \end{align}

            \STATE $\triangleright$ \textbf{Option II:} inexact inner solve
            \begin{align}
            \label{eq:dca-inner-solver-step}
            \bx^{k+1}
            =
            \mathcal S_{\mathcal A_k}(\bx^k,\bv^k;\Theta_k),
            \end{align}
            where $\mathcal S_{\mathcal A_k}$ denotes the inner solver selected
            by $\mathcal A_k$, and $\Theta_k$ collects its parameters and
            stopping tolerance.
        \ENDFOR
        \STATE {\bfseries Output:} $\hat\bx$, chosen uniformly from
        $\{\bx^k\}_{k=0}^{K-1}$.
    \end{algorithmic}
\end{algorithm} 

We use the unified DCA summarized in Algorithm~\ref{alg: DCA}. At the
$k$-th iteration, DCA computes a subgradient $\bv^k \in \partial g(\bx^k)$
and solves a $\gamma$-strongly convex subproblem to obtain $\bx^{k+1}$
(Option I). Option II is an abstract inexact inner-solver interface:
\Cref{alg:inexact-dca-inner} gives one realization, in which the convex
subproblem is solved to value accuracy by an inner first-order method,
and other inner-loop algorithms can be used as long as their output
satisfies the inexact certificate required in the analysis below.

Before analyzing the complexity of DCA, we first introduce the main
assumptions used throughout the section.

\subsection{Main Assumption}
Assume that $F$ is bounded below on $\operatorname{dom}\phi$, and denote
\[
F^\star:=\inf\limits_x F(x)>-\infty,
\qquad
\Delta_0:=F(\bx^0)-F^\star > 0.
\]
For any $\gamma>0$, $F$ admits the reformulation
\[
F = \Big(\phi+\tfrac{\gamma}{2}\|\bx\|^2\Big) - \Big(\tfrac{\gamma}{2}\|\bx\|^2+g\Big),
\]
so WLOG we take $\phi$ and $g$ to be $\gamma$-strongly convex, formalized as follows.

\begin{assumption}
\label{ass:strong-convexity-dc}
The functions $\phi$ and $g$ are strongly convex for some constant
$\gamma>0$, and $g$ has bounded subgradients:
$\sup\limits_{\bx}\sup\limits_{\bv\in\partial g(\bx)}\|\bv\|\le G$.
\end{assumption}

% We also take $\Delta_0>0$ WLOG: if $\Delta_0=0$, then $\bx^0$ is a global
% minimizer of $F$, so for any $\bv\in\partial g(\bx^0)$ and all $\bx$,
% \[
%     \phi(\bx)
%     \;\ge\;
%     g(\bx)+F(\bx^0)
%     \;\ge\;
%     g(\bx^0)+\langle\bv,\bx-\bx^0\rangle+F(\bx^0)
%     \;=\;
%     \phi(\bx^0)+\langle\bv,\bx-\bx^0\rangle,
% \]
% i.e., $\bv\in\partial\phi(\bx^0)\cap\partial g(\bx^0)$, so $\bx^0$ is already
% $(0,0,0)$-critical and the analysis is vacuous.
Our goal is to characterize the iteration complexity of DCA-type methods, without imposing smoothness on $g$, for reaching a $(\delta_\phi,\delta_g,\varepsilon)$-critical
point. As discussed below, the three parameters play distinct roles in this analysis:
$\varepsilon$ controls the residual mismatch between the approximate dual
vectors of $\phi$ and $g$, while $\delta_\phi$ and $\delta_g$ control how
loosely each convex component's Fenchel--Young gap is allowed to vanish.

\subsection{Common Auxiliary Lemmas}
% In particular, $\delta_\phi$ coincides with the primal optimality gap of the
% convex subproblem solved at each DCA iteration,
% \[
%     \min\limits_{\bx}\ \phi(\bx)-\langle \bv^k,\bx\rangle,
% \]
% evaluated at the returned iterate $\bx^{k+1}$. If the subproblem is solved to
% global optimality (an exact oracle), then $\Phi_\phi(\bx^{k+1},\bv^k)=0$,
% i.e., $\delta_\phi=0$; if it is solved only approximately -- e.g., by a fixed
% budget of inner first-order steps -- then $\delta_\phi>0$, and its magnitude
% is precisely the residual primal gap left by the inner solver. Exact
% subproblem solution therefore recovers the stronger regime $\delta_\phi=0$,
% in which $\delta_g$ remains the only source of inexactness, reflecting its
% structural rather than computational origin.

The next lemma is the pivot of the analysis, which shows that the $\phi$-side Fenchel gap $\Phi_\phi$ is at any candidate point equals the value suboptimality of that
point for the current convex subproblem. Consequently, the parameter
$\delta_\phi$ of \Cref{def:delta-eps-critical} coincides, exactly and
without slack, with the accuracy to which the subproblem is solved.

\begin{lemma}
\label{lem:dca-subproblem-value}
Let \Cref{ass:strong-convexity-dc} hold, $\bar\bx\in\dom\phi$,
$\bv\in\partial g(\bar\bx)$, and define
$P(\bx):=\phi(\bx)-\langle\bv,\bx\rangle$. Then the following statements
hold.
\begin{enumerate}
    \item[(i)] $\inf\limits_{\bx\in\RR^n}P(\bx)\ge
    F^\ast+g(\bar\bx)-\langle\bv,\bar\bx\rangle>-\infty$, and
    $\phi^*(\bv)=-\inf\limits_{\bx\in\RR^n}P(\bx)\in\RR$.
    \item[(ii)] For every $\by\in\dom\phi$, the following holds
    \[
        \Phi_\phi(\by,\bv)
        \;=\;
        P(\by)-\inf\limits_{\bx\in\RR^n}P(\bx).
    \]
    \item[(iii)] For every $\sigma\ge0$ and every $\by\in\dom\phi$, the
    following three statements are equivalent:
    \[
        \text{(a) }P(\by)\le\inf\limits_{\bx\in\RR^n}P(\bx)+\sigma;
        \qquad
        \text{(b) }\Phi_\phi(\by,\bv)\le\sigma;
        \qquad
        \text{(c) }\bv\in\partial_\sigma\phi(\by).
    \]
    In particular, $\by\in\argmin\limits_{\bx\in\RR^n}P(\bx)$ if and only if
    $\bv\in\partial\phi(\by)$.
\end{enumerate}
\end{lemma}

\begin{proof}
(i) Since $\bar\bx\in\dom\phi\subseteq\dom g$ and
$\bv\in\partial g(\bar\bx)$, the subgradient inequality gives
\begin{equation}
    g(\bx)\;\ge\;g(\bar\bx)+\langle\bv,\bx-\bar\bx\rangle
    \qquad\text{for all }\bx\in\RR^n,
    \label{eq:g-subgrad-at-xbar}
\end{equation}
equivalently
$\langle\bv,\bx\rangle\le g(\bx)-g(\bar\bx)+\langle\bv,\bar\bx\rangle$
for all $\bx\in\RR^n$. Hence, for every $\bx\in\dom\phi\subseteq\dom g$,
\[
    P(\bx)
    \;=\;
    \phi(\bx)-\langle\bv,\bx\rangle
    \;\ge\;
    \phi(\bx)-g(\bx)+g(\bar\bx)-\langle\bv,\bar\bx\rangle
    \;=\;
    F(\bx)+g(\bar\bx)-\langle\bv,\bar\bx\rangle
    \;\ge\;
    F^\ast+g(\bar\bx)-\langle\bv,\bar\bx\rangle,
\]
while $P(\bx)=+\infty$ for $\bx\notin\dom\phi$. Since $g(\bar\bx)\in\RR$,
the claimed lower bound follows, and it is finite. Furthermore,
$P(\bar\bx)=\phi(\bar\bx)-\langle\bv,\bar\bx\rangle<+\infty$, so
$\inf\limits P\in\RR$. By the definition of the Fenchel conjugate,
\[
    \phi^*(\bv)
    \;=\;
    \sup\limits_{\bx\in\RR^n}\bigl\{\langle\bv,\bx\rangle-\phi(\bx)\bigr\}
    \;=\;
    -\inf\limits_{\bx\in\RR^n}\bigl\{\phi(\bx)-\langle\bv,\bx\rangle\bigr\}
    \;=\;
    -\inf\limits_{\bx\in\RR^n}P(\bx)\;\in\;\RR .
\]

(ii) For $\by\in\dom\phi$, using (i),
\[
    \Phi_\phi(\by,\bv)
    \;=\;
    \phi(\by)+\phi^*(\bv)-\langle\bv,\by\rangle
    \;=\;
    P(\by)+\phi^*(\bv)
    \;=\;
    P(\by)-\inf\limits_{\bx\in\RR^n}P(\bx).
\]

(iii) The equivalence of (a) and (b) is immediate from (ii). The
equivalence of (b) and (c) is \Cref{lem:fenchel-gap-toolkit} applied to
$h=\phi$. The final claim is the case $\sigma=0$.
\end{proof}


The following identity decomposes the per-iteration decrease of $F$
into three nonnegative quantities that together constitute a
$(\delta_\phi,\delta_g,\eps)$-point. It holds for an arbitrary
candidate $\bx^{k+1}$, regardless of how it is produced, and therefore
serves as the common engine underlying both the exact and the
inexact analyses.

\begin{lemma}[One-step decomposition identity]
\label{lem:dca-one-step-identity}
Suppose \Cref{ass:strong-convexity-dc} holds. For any
$\bx^k,\bx^{k+1}\in\dom\phi$ and any $\bv^k\in\partial g(\bx^k)$, it holds that
\begin{equation}
    F(\bx^{k+1})-F(\bx^k)
    \;=\;
    \Phi_\phi(\bx^{k+1},\bv^k)
    \;-\;
    \Phi_\phi(\bx^k,\bv^k)
    \;-\;
    \Phi_g(\bx^{k+1},\bv^k).
    \label{eq:dca-one-step-identity}
\end{equation}
% where
% \[
%     \Phi_\phi(\bx^k,\bv^k)\ge0,
%     \qquad
%     \Phi_\phi(\bx^{k+1},\bv^k)\ge0,
%     \qquad
%     \Phi_g(\bx^{k+1},\bv^k)\ge0.
% \]
\end{lemma}

\begin{proof}
By the Fenchel--Young equality, i.e.,  for any 
\begin{equation}
    g^*(\bv^k)
    \;=\;
    \langle\bv^k,\bx^k\rangle-g(\bx^k)
    \;\in\;\RR,
    \label{eq:fenchel-young-equality-g}
\end{equation} and 
ssubstituting \eqref{eq:fenchel-young-equality-g} into the definition of
$\Phi_g$ gives
\begin{equation}
    \Phi_g(\bx^{k+1},\bv^k)
    \;=\;
    g(\bx^{k+1})+g^*(\bv^k)-\langle\bv^k,\bx^{k+1}\rangle
    \;=\;
    g(\bx^{k+1})-g(\bx^k)-\langle\bv^k,\bx^{k+1}-\bx^k\rangle.
    \label{eq:g-gap-bregman-form}
\end{equation}

On the other hand, rearranging the terms in the definition of
$\Phi_\phi$ yields
\[
    \Phi_\phi(\bx^k,\bv^k)-\Phi_\phi(\bx^{k+1},\bv^k)
    \;=\;
    \bigl[\phi(\bx^k)-\langle\bv^k,\bx^k\rangle\bigr]
    -\bigl[\phi(\bx^{k+1})-\langle\bv^k,\bx^{k+1}\rangle\bigr].
\]

Adding \eqref{eq:g-gap-bregman-form} to the above equality yields
\[
    \Phi_\phi(\bx^k,\bv^k)-\Phi_\phi(\bx^{k+1},\bv^k)
    +\Phi_g(\bx^{k+1},\bv^k)
    \;=\;
    \bigl[\phi(\bx^k)-g(\bx^k)\bigr]
    -\bigl[\phi(\bx^{k+1})-g(\bx^{k+1})\bigr]
    \;=\;
    F(\bx^k)-F(\bx^{k+1}),
\]
which completes the proof.
\end{proof}
By
\Cref{lem:dca-subproblem-value}\,(ii)--(iii), the quantity
$\Phi_\phi(\bx^{k+1},\bv^k)$ is the value inaccuracy of the subproblem
solve and populates the $\delta_\phi$-slot of the certificate at
$\bx^{k+1}$; the quantity $\Phi_g(\bx^{k+1},\bv^k)$ populates the
$\delta_g$-slot at $\bx^{k+1}$; and the quantity $\Phi_\phi(\bx^k,\bv^k)$
populates the $\delta_\phi$-slot of a second certificate anchored at
$\bx^k$, where the $g$-side is exact because $\Phi_g(\bx^k,\bv^k)=0$. The
objective decrease per iteration equals the sum of the two certificate
budgets minus the injected inexactness. In the exact regime the injected
inexactness vanishes and the identity strengthens the classical monotone
descent property of DCA \citep{lethi2009dc,abbaszadehpeivasti2023rate}
into an exact accounting of where the descent goes.

\subsection{Iteration Complexity of DCA-type Methods}
We translate the one-step identity above into finite-time complexity
bounds. We first treat exact DCA, where each convex subproblem is solved
exactly. We then analyze inexact DCA by separating the outer certificate
bound from the inner first-order cost.

\paragraph{Complexity analysis of Exact DCA.} We first analyze the complexity of DCA in \Cref{alg: DCA} with \emph{Option I}, meaning that each subproblem is solved exactly. By \Cref{lem:dca-subproblem-value}\,(ii)--(iii), the update
\eqref{eq:dca-exact-step} is equivalent to
\begin{equation}
    \Phi_\phi(\bx^{k+1},\bv^k)=0,
    \qquad\text{that is,}\qquad
    \bv^k\in\partial\phi(\bx^{k+1}).
    \label{eq:dca-exact-optimality}
\end{equation}

\begin{theorem}[Iteration Complexity of Exact DCA]
\label{thm:dca-exact-complexity}
Suppose \Cref{ass:strong-convexity-dc} holds. Let $\bx^0\in\dom\phi$ and
$\{\bx^k\}_{k\ge0}$ be generated by \Cref{alg: DCA} with \emph{Option I}.
Given $\delta_g>0$, DCA returns a point $\bar{\bx}\in\{\bx^1,\ldots,\bx^K\}$
satisfying $(0,\delta_g,0)$-critical within
$K=\mathcal O(\Delta_0 \delta_g^{-1})$ iterations.
\end{theorem}

\begin{proof}
Substituting the exact optimality condition~\eqref{eq:dca-exact-optimality}
into the one-step identity~\eqref{eq:dca-one-step-identity} gives, for
every $k\ge0$,
\begin{equation}
    F(\bx^k)-F(\bx^{k+1})
    \;=\;
    \Phi_\phi(\bx^k,\bv^k)+\Phi_g(\bx^{k+1},\bv^k)
    \;\ge\;0,
    \label{eq:dca-exact-descent}
\end{equation}
Summing~\eqref{eq:dca-exact-descent} over $k=0,\ldots,K-1$
and using $F(\bx^K)\ge F^\star$, we obtain
\[
    \frac1K\sum_{k=0}^{K-1}
    \Bigl[
        \Phi_\phi(\bx^k,\bv^k)+\Phi_g(\bx^{k+1},\bv^k)
    \Bigr]
    \;\le\;
    \frac{\Delta_0}{K}.
\] where defining
\[
    k^*:=\argmin_{0\le k\le K-1}
    \Bigl[
        \Phi_\phi(\bx^k,\bv^k)+\Phi_g(\bx^{k+1},\bv^k)
    \Bigr]
\]
yields
\[
    \Phi_\phi(\bx^{k^*},\bv^{k^*})+\Phi_g(\bx^{k^*+1},\bv^{k^*})
    \;\le\;\frac{\Delta_0}{K}.
\]
Because both summands are nonnegative, each is individually bounded by
their sum, implying
\begin{equation}
    \Phi_g(\bx^{k^*+1},\bv^{k^*})\;\le\;\frac{\Delta_0}{K}.
    \label{eq:dca-exact-individual-bound}
\end{equation}
Then, we set $\bu=\bv:=\bv^{k^*}$. Combining
\eqref{eq:dca-exact-optimality} with the optimality condition that $\bv^{k^*}\in\partial\phi(\bx^{k^*+1})\cap\partial g(\bx^k)$, we obtain $\bu\in\partial\phi(\bx^{k^*+1})$. By
\eqref{eq:dca-exact-individual-bound} together with
\Cref{lem:fenchel-gap-toolkit}, $\bv\in\partial_{\Delta_0/K}g(\bx^{k^*+1})$.
Hence the following holds
\[
    \operatorname{dist}
    \Bigl(
        \partial_{0}\phi(\bx^{k^*+1}),\,
        \partial_{\Delta_0/K}\,g(\bx^{k^*+1})
    \Bigr)
    \;\le\;
    \|\bu-\bv\|
    \;=\;
    0,
\] i.e., $\bx^{k^*+1}$ is $(0,\Delta_0/K,0)$-critical. Choosing $K=\lceil\Delta_0/\delta_g\rceil$ gives $\Delta_0/K\le\delta_g$.
Taking $\bar\bx:=\bx^{k^*+1}$ completes the proof.
\end{proof}

\blue{
\begin{remark}[Tightness under smooth oracle]
\label{rem:tightness-dimension}
\begin{itemize}
    \item \emph{Matching lower bound.} The rate of
    \Cref{thm:dca-exact-complexity} is tight up to the condition number
    $L/\gamma$: embedding the smooth hard instance of
    \citet{carmon2020lower} as the DC pair
    $\phi=\hat F+\tfrac{L}{2}\|\cdot\|^2$ and
    $g=\tfrac{L}{2}\|\cdot\|^2$ renders the component-wise DC oracle
    informationally equivalent to a standard first-order oracle, and
    yields a lower bound of
    $\Omega(\gamma L^{-1}\Delta\delta^{-1})$ oracle calls for reaching a
    $(0,\delta_g,\varepsilon)$-critical point with 
    $\delta_g=O(\gamma\varepsilon^2/L^2)$;  see \Cref{Appendix: lower bound of DCA}.

    % \item \emph{Dimension.} The embedded instance inherits the zero-chain
    % dimension $d=\Theta(\Delta L\varepsilon^{-2})$ of
    % \citet{carmon2020lower}; the bound itself is dimension-free. In
    % contrast, the one-dimensional barrier construction of
    % \Cref{thm:dc-oracle-limitation-sketch} trades dimension for component
    % growth $G_K=O(KL)$.

    \item \emph{Relation to classical measures.} When $g$ is
    $L_g$-smooth, $\delta_g$ and the same-point residual $\varepsilon$
    are explicitly inter-convertible with constant $L_g^2/\gamma$, so
    $O(\Delta\delta^{-1})$ recovers the classical
    $O(\varepsilon^{-2})$ rate; when $g$ is nonsmooth this equivalence
    provably fails, as the classical target is finite-time unattainable
    (\Cref{thm:dc-oracle-limitation-sketch}).
    % see \Cref{rem:smooth-conversion}.
\end{itemize}
\end{remark}}


\paragraph{Complexity Analysis of Inexact DCA} 

\begin{algorithm}[tbp]
    \caption{\textbf{Inexact DCA}}
    \label{alg:inexact-dca-inner}
    \begin{algorithmic}[1]
        \STATE {\bfseries Input:} Starting point $\bx^0\in\dom\phi$,
        tolerances $\{\sigma_k\}_{k\ge0}$.
        \FOR{$k=0,1,2,\ldots$}
            \STATE Choose $\bv^k\in\partial g(\bx^k)$.
            \STATE Define the convex subproblem
            \[
                P_k(\bx):=\phi(\bx)-\langle\bv^k,\bx\rangle .
            \]
            \STATE Starting from $\by^0=\bx^k$, run an inner first-order
            method to find $\by^{J_k}$ satisfying
            \begin{equation}
                P_k(\by^{J_k})
                \le
                \inf_{\bx\in\RR^n}P_k(\bx)+\sigma_k .
                \label{eq:dca-inexact-value-step}
            \end{equation}
            \STATE Set $\bx^{k+1}:=\by^{J_k}$.
        \ENDFOR
        \STATE {\bfseries Output:} $\{\bx^k\}_{k\ge0}$.
    \end{algorithmic}
\end{algorithm}

We now analyze the complexity of DCA with inexact subproblem solves, as
formalized in \Cref{alg:inexact-dca-inner}. This setting models the
practical situation in which the convex subproblem is solved by an
iterative method and terminated early \citep{yao2022large}. When
$\sigma_k=0$, condition \eqref{eq:dca-inexact-value-step} reduces to
\eqref{eq:dca-exact-step}.

\begin{lemma}
\label{lem:dca-inexact-complexity}
Suppose \Cref{ass:strong-convexity-dc} holds. Let
$\{\bx^k\}_{k\ge0}$ be generated by DCA with tolerances
$\{\sigma_k\}_{k\ge0}\subset[0,+\infty)$. Then the following statements
hold.
\begin{enumerate}
    \item[(i)] For every $k\ge0$, it holds that
    \[
        \Phi_\phi(\bx^{k+1},\bv^k)\;\le\;\sigma_k,
        \qquad\text{equivalently,}\qquad
        \bv^k\in\partial_{\sigma_k}\phi(\bx^{k+1}).
    \]
    \item[(ii)] For every integer $K\ge1$, it holds that
    \begin{equation}
        \sum_{k=0}^{K-1}
        \Bigl[
            \Phi_\phi(\bx^k,\bv^k)
            +\Phi_g(\bx^{k+1},\bv^k)
        \Bigr]
        \;\le\;
        \Delta_0+\sum_{k=0}^{K-1}\sigma_k .
        \label{eq:dca-inexact-telescope}
    \end{equation} Moreover, if  $\sigma_k\equiv\sigma\ge0$ for all $k$, then there exists $k^*\in\{0,1,\ldots,K-1\}$ such that
    \begin{equation}
        \Phi_\phi(\bx^{k^*},\bv^{k^*})
        +\Phi_g(\bx^{k^*+1},\bv^{k^*})
        \;\le\;
        \frac{\Delta_0}{K}+\sigma,
        \label{eq:dca-inexact-best-index}
    \end{equation} i.e.,  $\bx^{k^*+1}$ is a
        $\bigl(\sigma,\ \Delta_0/K+\sigma,\ 0\bigr)$-critical point.
\end{enumerate}
\end{lemma}

\begin{proof}
(i) Fix $k\ge0$. By \Cref{alg:inexact-dca-inner}, the iterate
$\bx^{k+1}$ satisfies the value-inexactness condition
\eqref{eq:dca-inexact-value-step} for the subproblem
$P_k(\bx)=\phi(\bx)-\langle\bv^k,\bx\rangle$. Applying
\Cref{lem:dca-subproblem-value}\,(iii) with
$\bar\bx=\bx^k$, $\bv=\bv^k$, $\by=\bx^{k+1}$, and
$\sigma=\sigma_k$ gives
\[
    \Phi_\phi(\bx^{k+1},\bv^k)\le\sigma_k
    \qquad\Longleftrightarrow\qquad
    \bv^k\in\partial_{\sigma_k}\phi(\bx^{k+1}).
\]
This proves (i).

(ii) Substituting the bound in (i) into the one-step identity
\eqref{eq:dca-one-step-identity} yields, for every $k\ge0$,
\begin{equation}
    F(\bx^k)-F(\bx^{k+1})
    \;=\;
    \Phi_\phi(\bx^k,\bv^k)
    -\Phi_\phi(\bx^{k+1},\bv^k)
    +\Phi_g(\bx^{k+1},\bv^k)
    \;\ge\;
    \Phi_\phi(\bx^k,\bv^k)
    +\Phi_g(\bx^{k+1},\bv^k)
    -\sigma_k .
    \label{eq:dca-inexact-descent}
\end{equation}
Summing~\eqref{eq:dca-inexact-descent} over $k=0,\ldots,K-1$
and using $F(\bx^K)\ge F^\star$, we obtain
\[
    \sum_{k=0}^{K-1}
    \Bigl[
        \Phi_\phi(\bx^k,\bv^k)+\Phi_g(\bx^{k+1},\bv^k)
    \Bigr]
    \;\le\;
    F(\bx^0)-F(\bx^K)+\sum_{k=0}^{K-1}\sigma_k
    \;\le\;
    \Delta_0+\sum_{k=0}^{K-1}\sigma_k,
\]
which proves \eqref{eq:dca-inexact-telescope}.

If $\sigma_k\equiv\sigma$, then the right-hand side of
\eqref{eq:dca-inexact-telescope} equals $\Delta_0+K\sigma$. Define
\[
    k^*
    :=
    \argmin_{0\le k\le K-1}
    \Bigl[
        \Phi_\phi(\bx^k,\bv^k)+\Phi_g(\bx^{k+1},\bv^k)
    \Bigr].
\]
Then
\[
    \Phi_\phi(\bx^{k^*},\bv^{k^*})
    +\Phi_g(\bx^{k^*+1},\bv^{k^*})
    \;\le\;
    \frac{\Delta_0}{K}+\sigma,
\]
which is \eqref{eq:dca-inexact-best-index}. Since
$\Phi_\phi(\bx^{k^*},\bv^{k^*})\ge0$, the same bound gives
\[
    \Phi_g(\bx^{k^*+1},\bv^{k^*})
    \;\le\;
    \frac{\Delta_0}{K}+\sigma .
\]
Part (i) gives
$\bv^{k^*}\in\partial_\sigma\phi(\bx^{k^*+1})$, while
\Cref{lem:fenchel-gap-toolkit} gives
\[
    \bv^{k^*}\in
    \partial_{\Delta_0/K+\sigma}g(\bx^{k^*+1}).
\]
Taking $\bu=\bv=\bv^{k^*}$, we obtain
\[
    \operatorname{dist}
    \Bigl(
        \partial_{\sigma}\phi(\bx^{k^*+1}),\,
        \partial_{\Delta_0/K+\sigma}g(\bx^{k^*+1})
    \Bigr)
    \;\le\;
    \|\bu-\bv\|
    \;=\;
    0.
\]
Hence $\bx^{k^*+1}$ is
$\bigl(\sigma,\Delta_0/K+\sigma,0\bigr)$-critical, completing the proof.
\end{proof}

\begin{theorem}[First-order Complexity of Inexact DCA]
\label{cor:dca-inexact-target}
Suppose \Cref{ass:strong-convexity-dc} holds. Let
$\{\bx^k\}_{k\ge0}$ be generated by \Cref{alg:inexact-dca-inner} with the
constant tolerance
\[
    \sigma_k\equiv\sigma_\star
    :=
    \min\{\delta_\phi,\theta\delta_g\},
    \qquad
    \theta\in(0,1),
\]
where $\delta_\phi,\delta_g>0$. Then DCA returns a
$(\delta_\phi,\delta_g,0)$-critical point
$\bar\bx\in\{\bx^1,\ldots,\bx^K\}$ within $K = \left(\Delta_0(\delta_g-\sigma_\star)^{-1}\right)$ outer iterations. Moreover, if each inner subproblem can be solved to value accuracy
$\sigma_\star$ using at most $\mathcal N_{\rm in}(\sigma_\star)$
inner first-order oracle calls, then the total number of 
first-order oracle calls is at most $ \left(\Delta_0(\delta_g-\sigma_\star)^{-1}  \mathcal N_{\rm in}(\sigma_\star) \right)$.
\end{theorem}

\begin{proof}
Applying \Cref{lem:dca-inexact-complexity}\,(ii) with
$\sigma=\sigma_\star$ gives an index $k^*\in\{0,\ldots,K-1\}$ such that
$\bx^{k^*+1}$ is
$\bigl(\sigma_\star,\Delta_0/K+\sigma_\star,0\bigr)$-critical. Since
$\sigma_\star\le\delta_\phi$ and
$K\ge\Delta_0/(\delta_g-\sigma_\star)$, we have
$\Delta_0/K+\sigma_\star\le\delta_g$. Hence
$\bx^{k^*+1}$ is $(\delta_\phi,\delta_g,0)$-critical. The oracle bound
follows by multiplying the number of outer iterations by the worst-case
inner cost per subproblem.
\end{proof}

\begin{remark}
\label{rem:dca-inexact-total-complexity}
We now specialize the total first-order complexity bound in
\Cref{cor:dca-inexact-target} by instantiating the inexact inner solver
with several standard first-order methods:
\[
\begin{array}{c|c|c}
\text{subproblem structure}
& \mathcal N_{\rm in}(\sigma_\star)
& \text{total first-order complexity}
\\ \hline
\text{convex, subgradient}
& O(\sigma_\star^{-2})
& O\!\left(\dfrac{\Delta_0}{\delta_g}\,\sigma_\star^{-2}\right)
\\
\text{strongly convex, subgradient}
& O(\gamma^{-1}\sigma_\star^{-1})
& O\!\left(\dfrac{\Delta_0}{\delta_g}\,\gamma^{-1}\sigma_\star^{-1}\right)
\\
\text{smooth strongly convex, GD}
& O(\kappa\log(1/\sigma_\star))
& O\!\left(\dfrac{\Delta_0}{\delta_g}\,\kappa\log(1/\sigma_\star)\right)
\\
\text{smooth strongly convex, APG}
& O(\sqrt{\kappa}\log(1/\sigma_\star))
& O\!\left(\dfrac{\Delta_0}{\delta_g}\,\sqrt{\kappa}\log(1/\sigma_\star)\right)
\end{array}
\]
The inner-solver bounds in the second column are the classical
deterministic first-order complexity estimates for subgradient, gradient,
and accelerated gradient methods; see, e.g.,
\citet[Chs.~3--4]{nesterov2018lectures} and
\citet[Chs.~3--5]{beck2017first}.
Here $\gamma$ is the strong convexity modulus of the subproblem $P_k$,
$\kappa=L_P/\gamma$ is its condition number when $P_k$ is $L_P$-smooth,
and the subgradient rows additionally assume the usual bounded-subgradient
constants, absorbed into $O(\cdot)$.

We highlight the insight of this result about the structural and
computational contributions to the total oracle complexity:
\blue{the outer factor $O(\Delta_0/\delta_g)$ is invariant under any
choice of inner solver, while the structure of the subproblem affects
only the multiplicative inner factor, which improves from
$O(\sigma_\star^{-2})$ to $O(\sqrt\kappa\log(1/\sigma_\star))$.}
\blue{Taking $\delta_\phi\to0$, so that $\mathcal N_{\rm in}(\sigma_\star)$
diverges, leaves the outer factor unchanged: even an exact subproblem
solver only certifies a $(0,\delta_g,0)$-critical point; see \Cref{thm:dca-exact-complexity}. This residual
$\delta_g$ therefore originates from the DCA mechanism
itself, namely the mismatch between the subgradient
$\bv^k\in\partial g(\bx^k)$ fixed at the old iterate and the certificate
evaluated at the new iterate $\bx^{k+1}$, rather than from any deficiency
of the inner solver.} This reflects, at the level of oracle complexity,
the asymmetry of \Cref{def:delta-eps-critical}: $\delta_g$ is a
structural cost of the nonsmooth component $g$, whereas $\delta_\phi$ is
a computational cost that enters only through $\mathcal N_{\rm in}$.
\blue{This decomposition constitutes an algorithmic rationale for
adopting two-sided tolerances in \Cref{def:delta-eps-critical}: merging
$\delta_\phi$ and $\delta_g$ into a single parameter forces the inner
accuracy to match the outer structural tolerance, eliminating the
freedom to decouple the two costs that the schedule
$\sigma_k\equiv\sigma_\star=\min\{\delta_\phi,\theta\delta_g\}$ in
\Cref{cor:dca-inexact-target} exploits. We emphasize that this invariance
is established within the DCA linearization framework and the two-point
certificate structure of \Cref{def:delta-eps-critical}; whether it
persists as a lower bound over a broader class of first-order algorithms
for Problem~\eqref{prob:main} is left open.}
\end{remark}

\subsection{Iteration Complexity of Proximal DCA-type Methods}
\label{subsec:prox-subgrad-dca}
We now analyze proximal DCA-type updates. A proximal term shifts the
$\phi$-side certificate from $\bv^k$ to a transported vector of the
form $\bu^{k+1}=\bv^k-\gamma^{-1}(\bx^{k+1}-\bx^k)$; the resulting
step-proportional residual $r\|\bx^{k+1}-\bx^k\|$ is what activates the
mismatch slot $\eps$, whereas bare DCA, which has no proximal term,
produces no shift and activates only the Fenchel-gap slots
$\delta_\phi,\delta_g$. The two deterministic updates---the single-loop
proximal-subgradient DCA (\Cref{alg:ps-dca}) and the exact proximal DCA
as its special case---share this mechanism exactly, so we isolate their
common estimates as \Cref{cond:certified-prox-descent}, prove the
complexity bound once at that level, and then verify the condition for
each. The stochastic variant of PS-DCA follows the same template but
adds a first-order gradient-error term to the residual; it is treated
separately in \Cref{subsec:stoch-ps-dca}.
Throughout this subsection, for
$\bv^k\in\partial g(\bx^k)$ and $s^k:=\bx^{k+1}-\bx^k$, define the
transported Fenchel gap of $g$
\begin{equation}
    \delta_g^{k+1}
    :=
    \Phi_g(\bx^{k+1},\bv^k)
    =
    g(\bx^{k+1})-g(\bx^k)
    -\langle \bv^k,\bx^{k+1}-\bx^k\rangle .
    \label{eq:prox-dca-transport-g-gap}
\end{equation}
The equality uses $\bv^k\in\partial g(\bx^k)$ and the Fenchel--Young
identity. In particular,
\[
    \bv^k\in\partial_{\delta_g^{k+1}} g(\bx^{k+1})
\]
by \Cref{lem:fenchel-gap-toolkit}: the subgradient $\bv^k$ that is exact
at $\bx^k$ becomes a $\delta_g^{k+1}$-subgradient once transported to
$\bx^{k+1}$.

\begin{condition}
\label{cond:certified-prox-descent}
Given $\bx^0\in\dom\phi$ and constants $a>0$, $r>0$, the sequences
$\{\bx^{k+1}\}$ and $\{\bu^{k+1}\}$, together with a choice
$\bv^k\in\partial g(\bx^k)$, satisfy the following: for every $k$ there
exists $\sigma_k\ge0$ such that the $\phi$-side
certificate
\begin{align}
    \bu^{k+1}
    &\in
    \partial_{\sigma_k}\phi(\bx^{k+1}),
    \label{eq:abstract-prox-phi-cert}
    \\
    \|\bu^{k+1}-\bv^k\|
    &\le
    r\|\bx^{k+1}-\bx^k\|,
    \label{eq:abstract-prox-dual-residual}
\end{align}
and the descent estimate
\begin{equation}
    F(\bx^k)-F(\bx^{k+1})
    \ge
    \delta_g^{k+1}
    +a\|\bx^{k+1}-\bx^k\|^2
    -\sigma_k
    \label{eq:abstract-prox-descent}
\end{equation}
hold.
\end{condition}

The next lemma is the abstract accounting device behind the
deterministic results in this subsection. It separates the two effects
in the certificate: the $\phi$-side Fenchel tolerance and the
transported $g$-side Fenchel tolerance.

\begin{lemma}
\label{lem:abstract-prox-certificate-bound}
Suppose $F^\star>-\infty$ and let $\Delta_0:=F(\bx^0)-F^\star$. Let
$K\ge1$, and suppose \Cref{cond:certified-prox-descent} holds for
$k=0,\ldots,K-1$ with
\[
    0\le\sigma_k\le\bar\sigma .
\]
Define
\[
    \hat k
    \in
    \argmin_{0\le k\le K-1}
    \bigl\{
        \delta_g^{k+1}
        +a\|\bx^{k+1}-\bx^k\|^2
    \bigr\},
    \qquad
    \bar\bx:=\bx^{\hat k+1}.
\]
Then $\bar\bx$ satisfies
\begin{align}
    \delta_g^{\hat k+1}
    &\le
    \frac{\Delta_0}{K}+\bar\sigma,
    \label{eq:abstract-prox-deltag-bound}
    \\
    \|\bx^{\hat k+1}-\bx^{\hat k}\|
    &\le
    \sqrt{
        \frac{1}{a}
        \left(
            \frac{\Delta_0}{K}+\bar\sigma
        \right)
    },
    \label{eq:abstract-prox-step-bound}
\end{align}
and hence
\begin{equation}
    \operatorname{dist}
    \left(
        \partial_{\bar\sigma}\phi(\bar\bx),
        \partial_{\frac{\Delta_0}{K}+\bar\sigma}g(\bar\bx)
    \right)
    \le
    r
    \sqrt{
        \frac{1}{a}
        \left(
            \frac{\Delta_0}{K}+\bar\sigma
        \right)
    } .
    \label{eq:abstract-prox-critical-bound}
\end{equation}
Consequently, $\bar\bx$ is a
\[
    \left(
        \bar\sigma,\,
        \frac{\Delta_0}{K}+\bar\sigma,\,
        r
        \sqrt{
            \frac{1}{a}
            \left(
                \frac{\Delta_0}{K}+\bar\sigma
            \right)
        }
    \right)\text{-critical point}.
\]
\end{lemma}

\begin{proof}
Summing the descent estimate \eqref{eq:abstract-prox-descent} over
$k=0,\ldots,K-1$ gives
\[
    \sum_{k=0}^{K-1}
    \left[
        \delta_g^{k+1}
        +
        a\|\bx^{k+1}-\bx^k\|^2
    \right]
    \le
    F(\bx^0)-F(\bx^K)+\sum_{k=0}^{K-1}\sigma_k .
\]
Since $F(\bx^K)\ge F^\star$ and $\sigma_k\le\bar\sigma$, we obtain
\begin{equation}
    \sum_{k=0}^{K-1}
    \left[
        \delta_g^{k+1}
        +
        a\|\bx^{k+1}-\bx^k\|^2
    \right]
    \le
    \Delta_0+K\bar\sigma .
    \label{eq:abstract-prox-telescope}
\end{equation}
By the definition of $\hat k$,
\[
    \delta_g^{\hat k+1}
    +
    a\|\bx^{\hat k+1}-\bx^{\hat k}\|^2
    \le
    \frac{\Delta_0}{K}+\bar\sigma .
\]
Both terms on the left-hand side are nonnegative, and hence
\eqref{eq:abstract-prox-deltag-bound} and
\eqref{eq:abstract-prox-step-bound} follow.

It remains to verify the same-point certificate. By
\eqref{eq:abstract-prox-phi-cert} and $\sigma_{\hat k}\le\bar\sigma$,
\[
    \bu^{\hat k+1}
    \in
    \partial_{\sigma_{\hat k}}\phi(\bx^{\hat k+1})
    \subseteq
    \partial_{\bar\sigma}\phi(\bx^{\hat k+1}).
\]
By the definition of $\delta_g^{\hat k+1}$ and
\eqref{eq:abstract-prox-deltag-bound},
\[
    \bv^{\hat k}
    \in
    \partial_{\delta_g^{\hat k+1}}g(\bx^{\hat k+1})
    \subseteq
    \partial_{\frac{\Delta_0}{K}+\bar\sigma}g(\bx^{\hat k+1}).
\]
Therefore
\[
\begin{aligned}
    \operatorname{dist}
    \left(
        \partial_{\bar\sigma}\phi(\bx^{\hat k+1}),
        \partial_{\frac{\Delta_0}{K}+\bar\sigma}g(\bx^{\hat k+1})
    \right)
    &\le
    \|\bu^{\hat k+1}-\bv^{\hat k}\|
    \\
    &\le
    r\|\bx^{\hat k+1}-\bx^{\hat k}\| .
\end{aligned}
\]
Combining this inequality with \eqref{eq:abstract-prox-step-bound}
proves \eqref{eq:abstract-prox-critical-bound}.
\end{proof}

The selection of $\hat k$ uses only computable quantities:
$\delta_g^{k+1}$ is evaluated from function values of $g$ and the
chosen $\bv^k$, and $\|\bx^{k+1}-\bx^k\|$ is observed, so $\hat k$ is
implementable without knowledge of the certificate tolerance
$\sigma_k$.

\begin{theorem}
\label{thm:abstract-prox-target}
Suppose the assumptions of \Cref{lem:abstract-prox-certificate-bound}
hold. Let target tolerances
$\delta_\phi,\delta_g,\eps>0$ be given, and choose
$\bar\sigma$ such that
\begin{equation}
    0\le\bar\sigma\le\delta_\phi,
    \qquad
    \bar\sigma
    <
    \min\left\{
        \delta_g,\,
        \frac{a\eps^2}{r^2}
    \right\}.
    \label{eq:abstract-prox-target-compatibility}
\end{equation}
If
\begin{equation}
    K
    \ge
    \left\lceil
    \Delta_0
    \max\left\{
        \frac{1}{\delta_g-\bar\sigma},\,
        \frac{1}{
            \frac{a\eps^2}{r^2}
            -
            \bar\sigma
        }
    \right\}
    \right\rceil ,
    \label{eq:abstract-prox-target-K}
\end{equation}
then the output $\bar\bx$ of
\Cref{lem:abstract-prox-certificate-bound} is a
$(\delta_\phi,\delta_g,\eps)$-critical point.
\end{theorem}

\begin{proof}
The first lower bound in \eqref{eq:abstract-prox-target-K} gives
\[
    \frac{\Delta_0}{K}+\bar\sigma
    \le
    \delta_g.
\]
The second lower bound gives
\[
    \frac{\Delta_0}{K}+\bar\sigma
    \le
    \frac{a\eps^2}{r^2}.
\]
Substituting these two inequalities into
\eqref{eq:abstract-prox-critical-bound} yields
\[
    \operatorname{dist}
    \left(
        \partial_{\bar\sigma}\phi(\bar\bx),
        \partial_{\delta_g}g(\bar\bx)
    \right)
    \le
    r\sqrt{
        \frac{1}{a}
        \cdot
        \frac{a\eps^2}{r^2}
    }
    =
    \eps .
\]
Finally, $\bar\sigma\le\delta_\phi$ implies
$\partial_{\bar\sigma}\phi(\bar\bx)
\subseteq
\partial_{\delta_\phi}\phi(\bar\bx)$, so $\bar\bx$ is
$(\delta_\phi,\delta_g,\eps)$-critical.
\end{proof}

We now verify \Cref{cond:certified-prox-descent} for concrete updates. We consider the composite case
\begin{equation}
    F(\bx)=f(\bx)+h(\bx)-g(\bx),
    \qquad
    \phi:=f+h,
    \label{eq:composite-ps-dca-model}
\end{equation}
where $f:\RR^n\to\RR$ is convex and $L_f$-smooth, and $h,g$ are
proper closed convex functions. We assume that the proximal mapping of
$h$ is easy to compute. This setting gives a single-loop realization of
the preceding certificate theorem.

\subsubsection{Proximal-Subgradient DCA}


For a stepsize $\eta\in(0,2/L_f)$, the proximal-subgradient DCA
(PS-DCA) update is
\begin{equation}
    \bx^{k+1}
    =
    \prox_{\eta h}
    \left(
        \bx^k-\eta(\nabla f(\bx^k)-\bv^k)
    \right),
    \qquad
    \bv^k\in\partial g(\bx^k).
    \label{eq:ps-dca-step}
\end{equation}
Equivalently,
\[
    \bx^{k+1}
    =
    \argmin_{\bx\in\RR^n}
    \left\{
        h(\bx)
        +
        \langle \nabla f(\bx^k)-\bv^k,\bx\rangle
        +
        \frac{1}{2\eta}\|\bx-\bx^k\|^2
    \right\}.
\]

Define
\[
    a_\eta:=\frac{1}{\eta}-\frac{L_f}{2}>0,
    \qquad
    r_\eta:=\frac{1}{\eta}+L_f .
\]

\begin{algorithm}[tbp]
    \caption{\textbf{Proximal-Subgradient DCA (PS-DCA)}}
    \label{alg:ps-dca}
    \begin{algorithmic}[1]
        \STATE {\bfseries Input:} $\bx^0\in\dom\phi$, stepsize
        $\eta\in(0,2/L_f)$, iteration number $K\ge1$.
        \FOR{$k=0,1,\ldots,K-1$}
            \STATE Choose $\bv^k\in\partial g(\bx^k)$.
            \STATE
            $\bx^{k+1}
            =\prox_{\eta h}
            \bigl(\bx^k-\eta(\nabla f(\bx^k)-\bv^k)\bigr)$.
        \ENDFOR
        \STATE Choose
        $\hat k\in\argmin_{0\le k\le K-1}
        \bigl\{\delta_g^{k+1}+a_\eta\|\bx^{k+1}-\bx^k\|^2\bigr\}$,
        with $\delta_g^{k+1}$ as in
        \eqref{eq:prox-dca-transport-g-gap}.
        \STATE {\bfseries Output:} $\bar\bx=\bx^{\hat k+1}$.
    \end{algorithmic}
\end{algorithm}

\begin{lemma}[PS-DCA realization]
\label{lem:ps-dca-realization}
Assume \eqref{eq:composite-ps-dca-model} and let
$\{\bx^k\}_{k\ge0}$ be generated by \eqref{eq:ps-dca-step}, with
$a_\eta,r_\eta$ defined above. Then \Cref{cond:certified-prox-descent}
holds with
\[
    a=a_\eta,
    \qquad
    r=r_\eta,
    \qquad
    \sigma_k=0 .
\]
More precisely, there exists
$\bu^{k+1}\in\partial\phi(\bx^{k+1})$ such that
\begin{align}
    \|\bu^{k+1}-\bv^k\|
    &\le
    r_\eta\|\bx^{k+1}-\bx^k\|,
    \label{eq:ps-dca-residual}
    \\
    F(\bx^k)-F(\bx^{k+1})
    &\ge
    \delta_g^{k+1}
    +
    a_\eta\|\bx^{k+1}-\bx^k\|^2 .
    \label{eq:ps-dca-descent}
\end{align}
\end{lemma}

\begin{proof}
Let $s^k:=\bx^{k+1}-\bx^k$. The optimality condition of
\eqref{eq:ps-dca-step} gives
\[
    \zero
    \in
    \partial h(\bx^{k+1})
    +
    \nabla f(\bx^k)
    -
    \bv^k
    +
    \frac{1}{\eta}s^k .
\]
Hence there exists $\bw^{k+1}\in\partial h(\bx^{k+1})$ such that
\[
    \bw^{k+1}
    =
    \bv^k
    -
    \nabla f(\bx^k)
    -
    \frac{1}{\eta}s^k .
\]
Since $f$ is convex and differentiable,
\[
    \bu^{k+1}
    :=
    \nabla f(\bx^{k+1})+\bw^{k+1}
    \in
    \partial(f+h)(\bx^{k+1})
    =
    \partial\phi(\bx^{k+1}).
\]
Moreover,
\[
\begin{aligned}
    \|\bu^{k+1}-\bv^k\|
    &=
    \left\|
        \nabla f(\bx^{k+1})-\nabla f(\bx^k)
        -
        \frac{1}{\eta}s^k
    \right\|
    \\
    &\le
    L_f\|s^k\|+\frac{1}{\eta}\|s^k\|
    =
    r_\eta\|s^k\|,
\end{aligned}
\]
which proves \eqref{eq:ps-dca-residual}.

It remains to prove the descent estimate. Define
\[
    Q_k(\bx)
    :=
    h(\bx)
    +
    \langle \nabla f(\bx^k)-\bv^k,\bx\rangle
    +
    \frac{1}{2\eta}\|\bx-\bx^k\|^2 .
\]
The function $Q_k$ is $\eta^{-1}$-strongly convex, and
$\bx^{k+1}$ minimizes $Q_k$. Hence
\[
    Q_k(\bx^k)-Q_k(\bx^{k+1})
    \ge
    \frac{1}{2\eta}\|s^k\|^2.
\]
Expanding this inequality gives
\begin{equation}
    h(\bx^k)-h(\bx^{k+1})
    +
    \langle \bv^k-\nabla f(\bx^k),s^k\rangle
    \ge
    \frac{1}{\eta}\|s^k\|^2 .
    \label{eq:ps-dca-h-descent}
\end{equation}
By $L_f$-smoothness of $f$,
\[
    f(\bx^{k+1})
    \le
    f(\bx^k)
    +
    \langle\nabla f(\bx^k),s^k\rangle
    +
    \frac{L_f}{2}\|s^k\|^2,
\]
or equivalently
\begin{equation}
    f(\bx^k)-f(\bx^{k+1})
    \ge
    -\langle\nabla f(\bx^k),s^k\rangle
    -
    \frac{L_f}{2}\|s^k\|^2 .
    \label{eq:ps-dca-f-descent}
\end{equation}
Combining \eqref{eq:ps-dca-h-descent} and
\eqref{eq:ps-dca-f-descent}, and using
\[
    \delta_g^{k+1}
    =
    g(\bx^{k+1})-g(\bx^k)-\langle\bv^k,s^k\rangle,
\]
we obtain
\[
\begin{aligned}
    F(\bx^k)-F(\bx^{k+1})
    &=
    f(\bx^k)-f(\bx^{k+1})
    +
    h(\bx^k)-h(\bx^{k+1})
    +
    g(\bx^{k+1})-g(\bx^k)
    \\
    &=
    f(\bx^k)-f(\bx^{k+1})
    +
    h(\bx^k)-h(\bx^{k+1})
    +
    \langle\bv^k,s^k\rangle
    +
    \delta_g^{k+1}
    \\
    &\ge
    \left(
        \frac{1}{\eta}
        -
        \frac{L_f}{2}
    \right)
    \|s^k\|^2
    +
    \delta_g^{k+1}.
\end{aligned}
\]
This proves \eqref{eq:ps-dca-descent}.
\end{proof}

\begin{corollary}[Iteration complexity of PS-DCA]
\label{cor:ps-dca-complexity}
Suppose $F^\star>-\infty$ and the assumptions of
\Cref{lem:ps-dca-realization} hold. Let
$\eta\in(0,2/L_f)$ and let $\{\bx^k\}_{k\ge0}$ be generated by
\eqref{eq:ps-dca-step}. Then for every $K\ge1$, there exists
$\hat k\in\{0,\ldots,K-1\}$ such that
$\bar\bx:=\bx^{\hat k+1}$ is a
\begin{equation}
    \left(
        0,\,
        \frac{\Delta_0}{K},\,
        r_\eta
        \sqrt{
            \frac{\Delta_0}{a_\eta K}
        }
    \right)\text{-critical point}.
    \label{eq:ps-dca-rate}
\end{equation}
Consequently, to obtain a $(0,\delta_g,\eps)$-critical point, it
suffices to take
\begin{equation}
    K
    \ge
    \left\lceil
    \max\left\{
        \frac{\Delta_0}{\delta_g},\,
        \frac{r_\eta^2\Delta_0}{a_\eta\eps^2}
    \right\}
    \right\rceil .
    \label{eq:ps-dca-target-K}
\end{equation}
In particular, for the standard choice $\eta=1/L_f$,
\[
    a_\eta=\frac{L_f}{2},
    \qquad
    r_\eta=2L_f,
\]
and \eqref{eq:ps-dca-target-K} becomes
\begin{equation}
    K
    \ge
    \left\lceil
    \max\left\{
        \frac{\Delta_0}{\delta_g},\,
        \frac{8L_f\Delta_0}{\eps^2}
    \right\}
    \right\rceil .
    \label{eq:ps-dca-standard-stepsize-K}
\end{equation}
The constant in the $\eps^{-2}$ term is minimized within the bound
\eqref{eq:ps-dca-target-K} by taking $\eta=1/(2L_f)$, which gives
\[
    \frac{r_\eta^2}{a_\eta}=6L_f .
\]
\end{corollary}

\begin{proof}
Apply \Cref{lem:abstract-prox-certificate-bound} with
$a=a_\eta$, $r=r_\eta$, and $\bar\sigma=0$, using
\Cref{lem:ps-dca-realization}. This gives
\[
    \delta_g^{\hat k+1}\le\frac{\Delta_0}{K},
    \qquad
    \|\bu^{\hat k+1}-\bv^{\hat k}\|
    \le
    r_\eta
    \sqrt{\frac{\Delta_0}{a_\eta K}},
\]
which is exactly \eqref{eq:ps-dca-rate}. The target bound
\eqref{eq:ps-dca-target-K} follows by requiring
$\Delta_0/K\le\delta_g$ and
$r_\eta\sqrt{\Delta_0/(a_\eta K)}\le\eps$.

For $\eta=1/L_f$,
\[
    \frac{r_\eta^2}{a_\eta}
    =
    \frac{(2L_f)^2}{L_f/2}
    =
    8L_f.
\]
For the last claim, set $a=1/\eta$ and minimize
\[
    \frac{(a+L_f)^2}{a-L_f/2}
    \qquad
    \text{over }a>L_f/2.
\]
The derivative vanishes at $a=2L_f$, i.e., at
$\eta=1/(2L_f)$, and the corresponding value is
\[
    \frac{(3L_f)^2}{3L_f/2}=6L_f.
\]
\end{proof}

\subsubsection{Exact Proximal DCA as a Special Case}

The classical exact proximal DCA update with parameter $\gamma>0$,
\begin{equation}
    \bx^{k+1}
    =
    \argmin_{\bx\in\RR^n}
    \left\{
        \phi(\bx)-\langle\bv^k,\bx\rangle
        +
        \frac{1}{2\gamma}\|\bx-\bx^k\|^2
    \right\},
    \qquad
    \bv^k\in\partial g(\bx^k),
    \label{eq:exact-prox-dca-step}
\end{equation}
is the instance of \eqref{eq:ps-dca-step} obtained by taking
$f\equiv0$ and $h=\phi$ with stepsize $\eta=\gamma$: the smooth part
vanishes, so $L_f=0$, every stepsize $\gamma>0$ is admissible, and the
update reduces to
$\bx^{k+1}=\prox_{\gamma\phi}(\bx^k+\gamma\bv^k)$, which is
\eqref{eq:exact-prox-dca-step}. The quadratic term makes the
subproblem strongly convex, so the minimizer is unique whenever
$\phi$ is proper closed convex; computing it is exactly the proximal
oracle for the full convex part $\phi$. \Cref{lem:ps-dca-realization}
applies with $L_f=0$ and yields the realization constants
\[
    a=\frac{1}{\gamma},
    \qquad
    r=\frac{1}{\gamma},
    \qquad
    \sigma_k=0,
\]
with the shifted dual certificate
$\bu^{k+1}=\bv^k-\gamma^{-1}(\bx^{k+1}-\bx^k)\in\partial\phi(\bx^{k+1})$.

\begin{corollary}[Iteration complexity of exact proximal DCA]
\label{cor:exact-prox-dca-complexity}
Suppose $F^\star>-\infty$, let $\phi$ and $g$ be proper closed convex
functions, and let $\{\bx^k\}_{k\ge0}$ be generated by
\eqref{eq:exact-prox-dca-step}. Then for every $K\ge1$, there exists
$\hat k\in\{0,\ldots,K-1\}$ such that
$\bar\bx:=\bx^{\hat k+1}$ is a
\begin{equation}
    \left(
        0,\,
        \frac{\Delta_0}{K},\,
        \sqrt{\frac{\Delta_0}{\gamma K}}
    \right)\text{-critical point}.
    \label{eq:exact-prox-dca-rate}
\end{equation}
Consequently, to obtain a $(0,\delta_g,\eps)$-critical point, it
suffices to take
\begin{equation}
    K
    \ge
    \left\lceil
    \max\left\{
        \frac{\Delta_0}{\delta_g},\,
        \frac{\Delta_0}{\gamma\eps^2}
    \right\}
    \right\rceil .
    \label{eq:exact-prox-dca-target-K}
\end{equation}
In particular, if $K\ge\lceil\Delta_0/(\gamma\eps^2)\rceil$, then the
same output is a $(0,\gamma\eps^2,\eps)$-critical point.
\end{corollary}

\begin{proof}
Apply \Cref{lem:abstract-prox-certificate-bound} with the constants
$a=r=\gamma^{-1}$ and $\bar\sigma=0$ supplied by
\Cref{lem:ps-dca-realization} with $f\equiv0$ and $\eta=\gamma$. Since
$r\sqrt{\Delta_0/(aK)}=\sqrt{\Delta_0/(\gamma K)}$, the abstract bound
\eqref{eq:abstract-prox-critical-bound} gives
\eqref{eq:exact-prox-dca-rate}, and
\eqref{eq:exact-prox-dca-target-K} follows by imposing
$\Delta_0/K\le\delta_g$ and $\sqrt{\Delta_0/(\gamma K)}\le\eps$.
\end{proof}

\begin{remark}[Realizations and activated slots]
\label{rem:prox-realization-table}
The two realizations above fit the same abstract certificate template:
\[
\begin{array}{c|c|c|c|c}
\text{method}
& a
& r
& \delta_\phi\text{-slot}
& \eps\text{-slot}
\\ \hline
\text{exact proximal DCA}
& \gamma^{-1}
& \gamma^{-1}
& 0
& \gamma^{-1}\|\bx^{k+1}-\bx^k\|
\\
\text{PS-DCA for }\phi=f+h
& \eta^{-1}-L_f/2
& \eta^{-1}+L_f
& 0
& (\eta^{-1}+L_f)\|\bx^{k+1}-\bx^k\|
\end{array}
\]
Thus exact proximal DCA and PS-DCA both activate the $\eps$-slot,
whereas classical bare DCA activates only the Fenchel-gap slots
$\delta_\phi,\delta_g$.
\end{remark}

\begin{remark}[Value-gap inexactness for the proximal subproblem]
\label{rem:prox-value-gap-warning}
The parameter $\sigma_k$ in
\eqref{eq:abstract-prox-phi-cert}--\eqref{eq:abstract-prox-descent} is a
certificate/decrease tolerance. A pure value-gap condition for the
proximalized subproblem,
\[
    M_k(\bx^{k+1})\le\inf_{\bx}M_k(\bx)+\sigma_k,
    \qquad
    M_k(\bx):=
    \phi(\bx)-\langle\bv^k,\bx\rangle
    +
    \frac{1}{2\gamma}\|\bx-\bx^k\|^2,
\]
does not by itself imply
\[
    \bv^k-\gamma^{-1}(\bx^{k+1}-\bx^k)
    \in
    \partial_{\sigma_k}\phi(\bx^{k+1}).
\]
Indeed, expanding the value-gap inequality around $\bx^{k+1}$ leaves an
additional term
\[
    -\frac{1}{2\gamma}\|\bz-\bx^{k+1}\|^2
\]
in the supporting inequality for $\phi$, which is unbounded below as a
function of $\bz$. For this reason, the unified theorem is stated in
terms of the verifiable certificate
\eqref{eq:abstract-prox-phi-cert} and the descent estimate
\eqref{eq:abstract-prox-descent}. Exact proximal DCA and PS-DCA satisfy
these conditions with $\sigma_k=0$.
\end{remark}

\begin{remark}[Coupling between $\delta_g$ and $\eps$]
\label{rem:prox-deltag-eps-coupling}
\Cref{cor:exact-prox-dca-complexity} gives the same output point with
\[
    \delta_g
    \le
    \frac{\Delta_0}{K},
    \qquad
    \eps
    \le
    \sqrt{\frac{\Delta_0}{\gamma K}}.
\]
Therefore, choosing $K=\lceil\Delta_0/(\gamma\eps^2)\rceil$ gives
\[
    \delta_g\le\gamma\eps^2.
\]
This is the precise proximal DCA coupling between the transported
Fenchel error and the dual certificate mismatch. It explains why the
proximal analysis is the natural source of an
$O(\Delta_0/(\gamma\eps^2))$ upper bound for the $\eps$-slot, while the
bare DCA analysis gives $\eps=0$ and only measures $\delta_g$.
\end{remark}

\subsection{Smoothing and the Two-Sided Certificate}
\label{sec:smoothing-certificate}

Smoothing the concave part is a third, and popular, route to nonsmooth DC
problems. Rather than propose a new method, we use it to make one point: the
$(\delta_{\phi},\delta_g,\eps)$-criticality of \Cref{def:delta-eps-critical} is
exactly the object such smoothing produces, so existing smoothing guarantees are
\emph{relabelings} of our certificate. Fix $\mu>0$ and recall the Moreau
envelope of $g$,
\begin{equation}
    g_\mu(\bx)
    :=
    \min_{\by\in\RR^n}
    \Bigl\{
        g(\by)+\tfrac{1}{2\mu}\|\by-\bx\|^2
    \Bigr\},
    \qquad
    \tilde{\bx}:=\prox_{\mu g}(\bx)
    =\argmin_{\by\in\RR^n}
    \Bigl\{
        g(\by)+\tfrac{1}{2\mu}\|\by-\bx\|^2
    \Bigr\},
    \label{eq:moreau-envelope}
\end{equation}
where the minimizer $\tilde{\bx}$ is unique because the inner objective is
$\tfrac1\mu$-strongly convex. Throughout this subsection we take
$(\phi,g)\in\mathcal D(\Delta,L,G)$, so both convex components are
finite-valued with uniformly bounded subgradients; accordingly we write
\[
    G_g:=\sup_{\bx}\;\sup_{\bv\in\partial g(\bx)}\|\bv\|\;\le\;G,
    \qquad
    G_\phi:=\sup_{\bx}\;\sup_{\bu\in\partial\phi(\bx)}\|\bu\|\;\le\;G .
\]
Replacing $g$ by $g_\mu$ renders the concave part smooth, and
$\eps$-criticality of the smoothed problem $\phi-g_\mu$ reads
$\operatorname{dist}\bigl(\partial\phi(\bx),\{\nabla g_\mu(\bx)\}\bigr)\le\eps$.
The next lemma shows that this smoothed certificate is, verbatim, a
$(\delta_{\phi},\delta_g,\eps)$-certificate for the original nonsmooth problem,
with an explicit $\delta_g$ that is linear in $\mu$.

\begin{lemma}[One-sided Moreau reduction]
\label{lem:moreau-reduction}
Let $(\phi,g)\in\mathcal D(\Delta,L,G)$, fix $\mu>0$, and suppose
$\bx\in\RR^n$ satisfies the $\eps$-criticality condition for the smoothed
problem $\phi-g_\mu$,
\[
    \operatorname{dist}\bigl(\partial\phi(\bx),\{\nabla g_\mu(\bx)\}\bigr)\le\eps ;
\]
equivalently, there exists $\bu\in\partial\phi(\bx)$ with
$\|\bu-\nabla g_\mu(\bx)\|\le\eps$. Then $\bx$ is a
$\bigl(0,\,2G_g^2\mu,\,\eps\bigr)$-critical point of the original problem
\eqref{prob:main} in the sense of
\Cref{def:delta-eps-critical}.
\end{lemma}

\begin{proof}
Write $\bv:=\nabla g_\mu(\bx)$ and $\tilde{\bx}:=\prox_{\mu g}(\bx)$.

\emph{Step 1 (Moreau identities).}
The optimality condition for $\tilde{\bx}$ in \eqref{eq:moreau-envelope} is
$\zero\in\partial g(\tilde{\bx})+\tfrac1\mu(\tilde{\bx}-\bx)$, hence
\begin{equation}
    \bv=\nabla g_\mu(\bx)=\frac{\bx-\tilde{\bx}}{\mu}\in\partial g(\tilde{\bx}).
    \label{eq:moreau-grad}
\end{equation}
Since $\bv\in\partial g(\tilde{\bx})$ and the subgradients of $g$ are bounded by
$G_g$, taking norms in \eqref{eq:moreau-grad} gives
\begin{equation}
    \|\bx-\tilde{\bx}\|=\mu\|\bv\|\le\mu G_g .
    \label{eq:prox-move}
\end{equation}

\emph{Step 2 (Transported Fenchel gap).}
Because $\bv\in\partial g(\tilde{\bx})$, the Fenchel--Young equality of
\Cref{lem:fenchel-gap-toolkit} at $\tilde{\bx}$ yields
$g^*(\bv)=\langle\bv,\tilde{\bx}\rangle-g(\tilde{\bx})$. Substituting this into
the gap function evaluated at the base point $\bx$,
\begin{equation}
    \Phi_g(\bx,\bv)
    =g(\bx)+g^*(\bv)-\langle\bv,\bx\rangle
    =g(\bx)-g(\tilde{\bx})-\langle\bv,\bx-\tilde{\bx}\rangle .
    \label{eq:gap-transported}
\end{equation}

\emph{Step 3 (Bounding the gap).}
Pick any $\bs_{\bx}\in\partial g(\bx)$, so $\|\bs_{\bx}\|\le G_g$; convexity of
$g$ gives the subgradient inequality
$g(\bx)-g(\tilde{\bx})\le\langle\bs_{\bx},\bx-\tilde{\bx}\rangle$. Combining this
with \eqref{eq:gap-transported}, the Cauchy--Schwarz and triangle inequalities,
\eqref{eq:moreau-grad}, and \eqref{eq:prox-move},
\begin{align}
    \Phi_g(\bx,\bv)
    &\le \langle\bs_{\bx}-\bv,\,\bx-\tilde{\bx}\rangle
    \le \|\bs_{\bx}-\bv\|\,\|\bx-\tilde{\bx}\|
    \le \bigl(\|\bs_{\bx}\|+\|\bv\|\bigr)\,\|\bx-\tilde{\bx}\| \notag\\
    &\le 2G_g\,\|\bx-\tilde{\bx}\|
    \le 2G_g^2\mu ,
    \label{eq:gap-bound}
\end{align}
where the last two inequalities use $\|\bs_{\bx}\|\le G_g$,
$\|\bv\|\le G_g$, and \eqref{eq:prox-move}.

\emph{Step 4 (Assembling the certificate).}
By \eqref{eq:gap-bound} and \Cref{lem:fenchel-gap-toolkit},
$\bv\in\partial_{2G_g^2\mu}\,g(\bx)$. The hypothesis furnishes
$\bu\in\partial\phi(\bx)=\partial_0\phi(\bx)$, so $\Phi_\phi(\bx,\bu)=0$, and
$\|\bu-\bv\|=\|\bu-\nabla g_\mu(\bx)\|\le\eps$. The pair $(\bu,\bv)$ therefore
witnesses
\[
    \Phi_\phi(\bx,\bu)\le 0,
    \qquad
    \Phi_g(\bx,\bv)\le 2G_g^2\mu,
    \qquad
    \|\bu-\bv\|\le\eps ,
\]
which is exactly the Fenchel-gap form of a $(0,2G_g^2\mu,\eps)$-critical point
in \Cref{def:delta-eps-critical}.
\end{proof}

\begin{corollary}[Two-sided DME reduction]
\label{cor:dme-reduction}
Let $(\phi,g)\in\mathcal D(\Delta,L,G)$ and fix $\mu_\phi,\mu_g>0$. Consider the
difference-of-Moreau-envelopes surrogate $\phi_{\mu_\phi}-g_{\mu_g}$, where
$\phi_{\mu_\phi}$ and $g_{\mu_g}$ are the Moreau envelopes of $\phi$ and $g$.
If $\bx$ is an $\eps$-critical point of this smooth surrogate, i.e.\
$\|\nabla\phi_{\mu_\phi}(\bx)-\nabla g_{\mu_g}(\bx)\|\le\eps$, then $\bx$ is a
$\bigl(2G_\phi^2\mu_\phi,\,2G_g^2\mu_g,\,\eps\bigr)$-critical point of the
original problem \eqref{prob:main}.
\end{corollary}

\begin{proof}
Apply the argument of \Cref{lem:moreau-reduction} to each component. Put
$\bu:=\nabla\phi_{\mu_\phi}(\bx)$ and $\bv:=\nabla g_{\mu_g}(\bx)$. Steps~1--3,
run on $g$ with parameter $\mu_g$, give
$\bv\in\partial_{2G_g^2\mu_g}\,g(\bx)$; the same steps, run symmetrically on
$\phi$ with parameter $\mu_\phi$ (with $\prox_{\mu_\phi\phi}$ in place of
$\prox_{\mu_g g}$ and $G_\phi$ in place of $G_g$), give
$\bu\in\partial_{2G_\phi^2\mu_\phi}\,\phi(\bx)$. The surrogate stationarity
$\|\bu-\bv\|\le\eps$ then closes the certificate through
\Cref{def:delta-eps-critical}.
\end{proof}

\begin{remark}[Smoothing relabels the target, not the rate]
\label{rem:smoothing-native-limit}
\Cref{lem:moreau-reduction} and \Cref{cor:dme-reduction} place existing
smoothing methods on the two-sided certificate and separate what their
guarantees keep from what they merely relabel. Driving $\mu\to0$ sends
$\delta_g=2G_g^2\mu\to0$, so smoothed stationarity recovers exact DC criticality
in the limit, while any fixed $\mu>0$ delivers a genuine same-point certificate
with an explicit, tunable $\delta_g=O(\mu)$.

\emph{What is unchanged---the iteration exponent.} In the deterministic regime,
the difference-of-Moreau-envelopes scheme of \citet{sun2023algorithms} reaches
an $\eps$-stationary point of the surrogate $\phi_{\mu_\phi}-g_{\mu_g}$ in
$O(\eps^{-2})$ proximal evaluations; composing with \Cref{cor:dme-reduction}
turns this into a $(2G_\phi^2\mu_\phi,2G_g^2\mu_g,\eps)$-certificate for the
original problem at the \emph{same} $O(\eps^{-2})$ rate reached here without any
smoothing (\Cref{cor:exact-prox-dca-complexity}). In the stochastic regime,
single-loop methods that optimize a smoothed objective---for instance
\citet{chayti2025stochastic} and \citet{hu2024single}---reach $\eps$-stationarity
in $O(\eps^{-4})$ iterations; since the reduction only relabels the point already
produced, this exponent is inherited unchanged on the two-sided certificate, the
smoothing parameter surfacing as the $\delta_g$ budget.

\emph{What changes---the target.} A single $\eps$-stationarity condition on a
smooth surrogate becomes the three-slot $(\delta_\phi,\delta_g,\eps)$-certificate
of the original nonsmooth problem, in which the smoothing parameters appear as
the explicit, separately controllable budgets $\delta_\phi=O(\mu_\phi)$ and
$\delta_g=O(\mu_g)$. The bias hidden inside the surrogate---the discrepancy
between the smoothed and the original objective---is thereby exposed as a
certificate slot rather than left implicit. Smoothing is consequently dispensable
in the deterministic regime, where the transported Fenchel gap $\delta_g$ already
absorbs the nonsmoothness of $g$ directly; it re-enters only as an algorithmic
device that renders a component differentiable, which is precisely the structure
a stochastic variance-reduced analysis can exploit.
\end{remark}

